\chapter{Algèbre qutrit de couleur et connexion finie de jauge}
\label{chap:algebra-color-qutrit}
\index{color!qutrit}\index{opérateurs de Weyl}
\index{groupe spécial unitaire}\index{action de Wilson}

La fonctionnelle résiduelle \(\kappa_3\) du
\cref{chap:ocupacion-estadistica} fournit trois classes affines ; l'atlas
ultérieur réalise leur bilan \(12+12+12\) dans la
\cref{prop:color-atlas}. Un triplet d'étiquettes n'est pas encore une
représentation. L'étape algébrique requiert de déclarer la carte qutrit, son
orientation et son caractère de phase. Sous ces hypothèses, la représentation et
son algèbre infinitésimale s'obtiennent exactement, sans utiliser de masses ni de données
chromodynamiques.

\section{Couple de Weyl et algèbre de trace \(0\)}
\label{sec:weyl-color-qutrit}

Soit
\[
 V_{\rm col}=\mathbb C^3,
 \qquad
 \omega=e^{2\pi i/3},
\]
avec la base ordonnée \((e_0,e_1,e_2)\). Nous définissons
\[
 Xe_j=e_{j+1\bmod3},
 \qquad
 Ze_j=\omega^j e_j.
\]
On a
\[
 X^3=Z^3=I,
 \qquad
 ZX=\omega XZ.
\]

\begin{theorem}[Base de Weyl et algèbre complexe de couleur]
\label{thm:weyl-sl3-color}
Pour
\[
 W_{ab}:=X^aZ^b,
 \qquad
 (a,b)\in(\mathbb Z/3\mathbb Z)^2,
\]
les neuf opérateurs \(W_{ab}\) forment une base orthogonale de
\(\operatorname{End}_{\mathbb C}(V_{\rm col})\) pour le produit de
Hilbert--Schmidt. Les huit opérateurs différents de l'identité ont une trace nulle et
\[
 \operatorname{span}_{\mathbb C}
 \{W_{ab}:(a,b)\ne(0,0)\}
 =\mathfrak{sl}_3(\mathbb C).
\]
\end{theorem}

\begin{proof}
La relation de Weyl permet de réduire tout produit à la forme
\(X^aZ^b\). Un calcul direct donne
\[
 \operatorname{tr}(W_{ab}^{\dagger}W_{cd})
 =3\,\delta_{ac}\delta_{bd}.
\]
Par conséquent, les neuf opérateurs sont linéairement indépendants ; comme
\(\dim_{\mathbb C}\operatorname{End}_{\mathbb C}(V_{\rm col})=9\), ils forment une
base. La trace de \(X^aZ^b\) s'annule sauf pour
\((a,b)=(0,0)\). Les huit autres forment ainsi une base du sous-espace de
trace nulle, dont la dimension complexe est huit.
\end{proof}

\begin{corollary}[Forme réelle compacte]
\label{cor:su3-color-qutrit}
L'involution
\[
 \tau(A):=-A^\dagger
\]
préserve \(\mathfrak{sl}_3(\mathbb C)\), et son ensemble de points fixes est
\[
 \mathfrak{sl}_3(\mathbb C)^\tau
 =\{A:A^\dagger=-A,\ \operatorname{tr}A=0\}
 =\mathfrak{su}(3).
\]
Cette forme réelle est de dimension huit.
\end{corollary}

\begin{proof}
On a
\[
 W_{ab}^{\dagger}=\omega^{ab}W_{-a,-b}.
\]
Les huit indices non nuls forment quatre couples adjoints. Pour
\[
 \mathcal R=\{(1,0),(0,1),(1,1),(1,2)\},
\]
les opérateurs
\[
 A_u=W_u-W_u^\dagger,
 \qquad
 B_u=i(W_u+W_u^\dagger),
 \qquad u\in\mathcal R,
\]
sont antihermitiens, de trace nulle et linéairement indépendants sur
\(\mathbb R\). Ils constituent donc une base réelle de
\(\mathfrak{su}(3)\).
\end{proof}

L'application résiduelle
\[
 \kappa_{\rm col}(r,c)
 =\mathbb C e_{\,r-c\bmod3}
\]
associe les trois classes de l'atlas aux trois droites propres de \(Z\).
L'opérateur \(X\) les permute cycliquement et \(Z\) enregistre leur phase. Cette
attribution est exacte une fois fixés
\[
 V_{\rm col}=\mathbb C[\mathbb F_3],
 \quad
 \text{la base delta},
 \quad
 j\longmapsto j+1,
 \quad
 j\longmapsto\omega^j.
\]
Les changements affines admissibles de la carte produisent des représentations
unitairement conjuguées. La construction est donc naturelle à
conjugaison unitaire près ; elle ne sélectionne pas à elle seule les noms physiques des
couleurs ni une réalisation QCD.

\section{Connexion et action sur un complexe cellulaire fini}
\label{sec:gauge-finito-color}

\begin{construction}[Théorie de jauge finie relative au complexe]
\label{cons:gauge-finito-color}
Soit \(\mathcal K=(V,E,F)\) une cellularisation finie orientée compatible avec
le graphe des chemins. À chaque arête orientée \(e\), on attribue
\[
 U_e\in SU(3),
 \qquad
 U_{\bar e}=U_e^{-1}.
\]
Une transformation \(g:V\to SU(3)\) agit par
\[
 U_e\longmapsto U_e^g
 =g_{t(e)}U_eg_{s(e)}^{-1}.
\]
Pour chaque face \(f\), le produit ordonné
\[
 U_f=\prod_{e\subset\partial f}U_e^{\epsilon(f,e)}
\]
définit son holonomie. Un sommet de base \(v_f\) étant choisi, on a
\[
 U_f^g=g_{v_f}U_fg_{v_f}^{-1}.
\]
\end{construction}

\begin{proposition}[Invariance de jauge et transport covariant]
\label{prop:wilson-color-finito}
Pour \(\beta\ge0\), l'action
\[
 S_{\rm col}^{(\mathcal K,\beta)}[U]
 =\frac{\beta}{3}\sum_{f\in F}
 \left(3-\operatorname{Re}\operatorname{tr}U_f\right)
\]
est invariante de jauge et non négative. Si
\(\psi_v\in V_{\rm col}\) se transforme comme
\(\psi_v\mapsto g_v\psi_v\), alors
\[
 \nabla_e\psi:=U_e\psi_{s(e)}-\psi_{t(e)}
\]
satisfait
\[
 \nabla_e^g\psi^g=g_{t(e)}\nabla_e\psi.
\]
\end{proposition}

\begin{proof}
L'holonomie change par conjugaison et la trace est cyclique ; d'où
l'invariance. Comme chaque \(U_f\) est unitaire de dimension trois,
\(\operatorname{Re}\operatorname{tr}U_f\le3\), ce qui prouve la
non-négativité. La covariance de \(\nabla_e\) résulte de la substitution des deux lois
de transformation.
\end{proof}

\section{Positivité des noyaux de classe et décomposition harmonique}

La connexion finie précédente n'implique pas à elle seule la positivité par
réflexion ni une limite de Yang--Mills. Avant de formuler de tels transferts,
il convient d'isoler le contrôle classique qui est effectivement clos. Soit \(G\) un groupe
compact et \(\widehat G\) son dual unitaire. Pour chaque
\(\rho\in\widehat G\), soient \(d_\rho\) sa dimension et
\(c_\rho(a)\ge0\) des coefficients de décroissance suffisante, par rapport à la
croissance de \(d_\rho\), pour que la série suivante converge
absolument.

\begin{definition}[Noyau de classe]
\label{pf84:YM-07}
On définit
\[
 K_a(g,h):=
 \sum_{\rho\in\widehat G}
 c_\rho(a)\,
 \Tr\!\bigl(\rho(g)\rho(h)^*\bigr),
 \qquad g,h\in G.
\]
Le choix
\[
 c_\rho(t)=d_\rho e^{-tC_2(\rho)},
 \qquad t>0,
\]
fournit l'exemple du noyau de chaleur. La non-négativité des
coefficients ne remplace pas l'hypothèse de convergence absolue.
\end{definition}

\begin{theorem}[Positivité de Peter--Weyl]
\label{pf84:YM-08} Pour \(g_1,\dots,g_n\in G\) et \(z_1,\dots,z_n\in\CC\),
\[
 \sum_{i,j=1}^n
 \overline z_i z_jK_a(g_i,g_j)\ge0.
\]
\end{theorem}

\begin{proof}
L'unitarité des représentations et la convergence absolue permettent
de réordonner la somme :
\[
 \sum_{i,j}\overline z_i z_jK_a(g_i,g_j)
 =
 \sum_{\rho\in\widehat G}
 c_\rho(a)
 \left\lVert
 \sum_j\overline z_j\rho(g_j)
 \right\rVert_{\mathrm{HS}}^2\ge0.
\]
Un coefficient spectral négatif peut détruire cette positivité de Gram ;
le signe fait donc partie des hypothèses.
\end{proof}

\begin{lemma}[Contrôle abélien de Poisson]
\label{pf84:YM-09}
Pour \(G=U(1)\), \(\theta\in\RR\) et \(|a|<1\),
\[
 \sum_{n\in\ZZ}a^{|n|}e^{in\theta}
 =\frac{1-a^2}{1-2a\cos\theta+a^2}.
\]
Si l'on exige en outre des coefficients de Fourier non négatifs, on restreint
\(a\) à \(0\le a<1\).
\end{lemma}

\begin{proof}
L'égalité résulte de la sommation séparée des deux séries géométriques
d'indices positifs et négatifs et de l'ajout du mode \(n=0\). Pour \(a<0\),
l'identité scalaire reste valable, mais ne constitue plus un exemple de
coefficients de Peter--Weyl non négatifs.
\end{proof}

\begin{remark}
Les \cref{pf84:YM-07,pf84:YM-08,pf84:YM-09} sont des définitions et des résultats
classiques d'analyse harmonique. Ils fournissent un contrôle positif pour
les noyaux de classe ; ils ne démontrent ni la positivité d'Osterwalder--Schrader, ni l'écart de
masse, ni le confinement, ni la limite continue, et ne promeuvent pas la théorie de jauge
finie de ce chapitre au rang de chromodynamique quantique.
\end{remark}

\begin{exactbox}[title={Clôture algébrique de la couleur}]
Les \cref{thm:weyl-sl3-color,cor:su3-color-qutrit} construisent exactement la
chaîne
\[
\text{reste ternaire}\longrightarrow
\text{qutrit}\longrightarrow
\mathfrak{sl}_3(\mathbb C)\longrightarrow\mathfrak{su}(3).
\]
La \cref{cons:gauge-finito-color,prop:wilson-color-finito} prolonge cette
chaîne en une connexion et une action de jauge sur la cellularisation finie
\(\mathcal K\). Tel est l'objet démontré dans le chapitre. La
comparaison avec la chromodynamique physique s'effectue ensuite et ne redéfinit ni le
qutrit, ni le bilan \(12+12+12\), ni la représentation construite.
\end{exactbox}
